\section{Introduction}
\label{sec:introduction}

Thomason's model structure equips the category of small categories with
the homotopy theory of spaces. Its higher-dimensional analogue asks for
the same result on the entire category of strict globular
$n$-categories, with weak equivalences detected by the Street nerve and
fibrations detected by its twice-extended nerve. We prove this analogue
in every dimension, resolving positively the higher-dimensional
Thomason model-structure conjecture of Ara and Maltsiniotis.

Fix a universe of small sets and enlarge the universe when forming
categories of small objects. For a positive integer $n$, let $\nCat$ be
the category of small strict globular $n$-categories and strict functors.
We regard it as the full subcategory of $\omCat$ consisting of strict
$\omega$-categories whose cells above dimension $n$ are identities; for
$n=\infty$, set $\nCat=\omCat$. There is no invertibility assumption on
positive-dimensional cells.
The Street nerve and its left adjoint are denoted by
\[
 c_n:\sSet\rightleftarrows\nCat:N_n,
 \qquad (N_nX)_m=\Hom_{\omCat}(\Orient{m},X),
\]
where $\Orient{m}$ is the $m$th oriental. With the usual subdivision
adjunction $\Sd\dashv\Ex$, put
\[
 L_n=c_n\Sd^2,\qquad R_n=\Ex^2N_n.
\]
Write $W_{\KQ}$ and $\mathrm{Fib}_{\KQ}$ for the weak homotopy
equivalences and Kan fibrations of simplicial sets, and define
\begin{equation}
\label{eq:prescribed-classes}
 W_n=R_n^{-1}(W_{\KQ})=N_n^{-1}(W_{\KQ}),\qquad
 F_n=R_n^{-1}(\mathrm{Fib}_{\KQ}),\qquad
 C_n={}^{\perp}(F_n\cap W_n).
\end{equation}
Here ${}^{\perp}$ denotes the left lifting class. The equality defining
$W_n$ follows from the classical properties of $\Ex$, recalled in
Section~\ref{sec:transfer}.

\begin{theorem}\label{thm:main}
For every $n\in\{1,2,3,\ldots\}\cup\{\infty\}$, the classes
$(C_n,W_n,F_n)$ in \eqref{eq:prescribed-classes} form a proper
combinatorial model structure on $\nCat$. It is cofibrantly generated by
$(L_nI,L_nJ)$, where
\[
 I=\{\partial\Delta[m]\hookrightarrow\Delta[m]:m\ge0\},\qquad
 J=\{\Lambda^k[m]\hookrightarrow\Delta[m]:m\ge1,\ 0\le k\le m\}.
\]
Moreover,
\[
 L_n:\sSet_{\KQ}\rightleftarrows(\nCat,C_n,W_n,F_n):R_n
\]
is a Quillen equivalence.
\end{theorem}

\subsection*{History and significance}
Thomason established the case $n=1$ by transferring the simplicial
model structure through $c_1\Sd^2\dashv\Ex^2N_1$
\cite{Thomason1980}. His proof separates two issues that persist in
higher dimensions: the homotopy behavior of the adjunction unit and
the comparison between a pushout of nerves and the nerve of a
categorical pushout. For the latter he introduced Dwyer maps and proved
the comparison for their arbitrary cobase changes
\cite[Definition~4.1 and Proposition~4.3]{Thomason1980}.
Cisinski corrected the retract argument used for properness, while
preserving the model structure and its properness
\cite{Cisinski1999}.
Street's orientals extend the nerve to strict higher categories
\cite{Street1987}, providing the weak equivalences in the question
considered here.

Ara and Maltsiniotis developed an abstract transfer theorem, drawing
on the homotopy theory of categories developed by Grothendieck and
Cisinski, and established the proper combinatorial model structure
in dimension two. Using Chiche's comparison theorem for the
homotopy categories of categories and $2$-categories, they also
obtained the Quillen equivalence
\cite[Theorems~4.11, 6.27, 6.30, and Corollary~6.32]{AraMaltsiniotis2014}.
Worytkiewicz, Hess, Parent, and Tonks had proposed the two-dimensional
structure earlier \cite{WorytkiewiczHessParentTonks2007}; their later
corrigendum addresses the gaps identified by Ara and Maltsiniotis
\cite{WorytkiewiczHessParentTonks2016}.
For arbitrary $1\le n\le\infty$, the transfer framework isolates two
sufficient conditions: a unit comparison on nerves of posets, and
preservation of certain poset equivalences by every categorical
pushout \cite[Scholie~5.14]{AraMaltsiniotis2014}.

Ara and Maltsiniotis proved the poset-unit condition in all these
dimensions \cite[Corollary~10.11 and Scholie~10.12]{AraMaltsiniotis2015}.
Gagna subsequently gave a direct simplicial proof and showed that
the Street nerve induces an equivalence with homotopy types after
localizing at the nerve weak equivalences
\cite[Theorems~5.6 and~6.9]{Gagna2018}.
Thus strict higher categories already described the homotopy category
of spaces in every dimension. The prescribed model structure requires
the additional lifting and factorization statements, whose construction
depends on controlling actual categorical pushouts. Ara's 2023 account
records this remaining sufficient condition and the general conjecture
\cite[\S\S9.2.5--9.2.6]{Ara2023}; Guetta and Maltsiniotis also state
the $\omega$-dimensional model-structure conjecture explicitly
\cite[Introduction, p.~6]{GuettaMaltsiniotis2024}.

The shapes of higher cells matter in this question. Fiore and Paoli
constructed a Thomason model structure on small $n$-fold categories,
whose weak equivalences are detected by the diagonal of the $n$-fold
nerve and which is Quillen equivalent to simplicial sets
\cite[Theorem~9.28]{FiorePaoli2010}.
Their objects are obtained by iterating the internal-category
construction; the present theorem concerns strict globular categories
and the Street nerve.

Theorem~\ref{thm:main} supplies both weak factorization systems on all
strict $n$-categories, together with left and right properness, by
establishing the remaining pushout condition. Its arbitrary-target
scope is essential: the small object argument must allow attachments
to categories that need not be cofibrant or freely generated.
The $\omega$-dimensional proof uses finite orientals directly, and
finite-dimensional reflection carries the same construction to each
finite $n$.

\subsection*{Stationary cylinders and the pushout argument}
For a poset $E$, write $\U{E}=c_n\OrdN E$, where $\OrdN$ is the ordinary
nerve. The principal intermediate result is
Theorem~\ref{thm:couniversal}: if $A\subseteq E$ is a downward-closed
full subposet and its inclusion has a right adjoint retraction, then
every pushout of $\U{A}\to\U{E}$ is a Street weak equivalence.
This proves the sufficient condition isolated in
\cite[Scholie~5.14, condition~(d$'$)]{AraMaltsiniotis2014}.
The cylinders used for contraction fix $\U{A}$ throughout the interval.
After any attachment $\U{A}\to X$, this equality lets each
cylinder glue to the constant cylinder on $X$, controlling the entire
pushout and its newly formed composites.

The construction starts with a restricted but explicit directed prism.
Suppose $f,g:E\to E$ are monotone, $f\le g$, they agree on a sieve $A$,
and $g$ is constant on $E\setminus A$. Write $u$ for the degree-one
generator of the interval chain complex. For every increasing tuple $w$,
the formula
\[
 [w_0,\ldots,w_p]\otimes u
 \longmapsto [fw_0,\ldots,fw_p,gw_p]
\]
together with the endpoint maps induced by $f$ and $g$ defines a positive
chain map after normalized repetitions are set to zero. It produces a
strict Gray cylinder constant in the interval direction on $\U{A}$,
at the common endpoint restriction induced by $f|_A=g|_A$.
Its positivity is a consequence of the precise hypotheses on the old
part and the constant upper map. Section~\ref{sec:prisms} proves the
chain identity, compatibility with all simplex operators, and descent
to every finite dimension.

When $E\setminus A$ has a greatest element, we choose two such
cylinders whose defining poset maps restrict to $\id_A$.
They give a relative contraction by a zigzag, which extends across
any attachment. For a general finite poset, deleting
distinct maximal new elements gives two sieves containing $A$.
The corresponding pushouts are full sieve subcategories of the total
attachment, with intersection the pushout for the common subposet.
Their nerves form an ordinary pushout along
monomorphisms, so the local weak-equivalence statements glue. This
allows each local contraction to be chosen independently. Finite
subposets stable under the retraction then yield the result for
arbitrary posets.

The remaining comparison follows the geometry of Thomason's Dwyer
maps \cite[Propositions~4.2--4.3]{Thomason1980}.
Two subdivisions put each generating boundary or horn inclusion into
the form $\OrdN A\hookrightarrow\OrdN B$, where $A$ is a sieve and
its inclusion into its upward closure in $B$ has a right adjoint
retraction (Lemma~\ref{lem:subdivision-posets}). For any small poset
pair with these properties, Proposition~\ref{prop:attachment-comparison}
provides a reusable comparison between the pushout of Street nerves
and the nerve of an attachment along an arbitrary functor $\U{A}\to X$.
The poset description here concerns the generating pieces.

At the chain level we use Steiner's description of orientals
\cite{Steiner2004,Steiner2007} and the biclosed Gray tensor product,
with the completed extension theorem of Ara and Maltsiniotis
\cite[Theorem~A.15]{AraMaltsiniotis2020}.
Positive-chain constructions already give simplicial nerve homotopies
in Gagna's proof \cite[Sections~6.3--6.8]{Gagna2018}.
Here the restricted prism hypotheses give strict relative cylinders;
the choices fixing $A$ extend across arbitrary attachments.
The cover calculations build on the
fiber and sieve machinery of Ara and Maltsiniotis
\cite[Lemma~2.8 and Propositions~5.3 and~5.5]{AraMaltsiniotis2014};
they allow us to glue the weak-equivalence conclusions without choosing
compatible contractions on overlapping pieces.

\subsection*{Organization}
Section~\ref{sec:foundations} records the foundations and the precise
simplicial gluing lemma. Sections~\ref{sec:prisms} and
\ref{sec:pieces} construct the stationary cylinders and identify
labelled pieces. Section~\ref{sec:couniversality} proves couniversality
and the unit comparison for posets. Section~\ref{sec:attachments}
proves the nerve comparison for subdivided boundary and horn
attachments, with explicit control of the monomorphic gluing legs.
Section~\ref{sec:transfer} supplies the model axioms, properness, and
the unit comparison for arbitrary simplicial sets.
Appendix~\ref{app:presentability} proves the required presentability
and smallness assertions by guarded term presentations.
Appendix~\ref{sec:context} discusses the relationship with the
elementary-expansion method for Grothendieck groupoids; its claims are not inputs to the
proof.
